% ATTEMPT_STATUS: unresolved
% RESEARCH_GENERATED: 2026-10-01; GPT-6 Astra Ultra; individual mathematical attempt
% TOP_PROBLEM: {"id": 139, "title": "Axis-Parallel Rectangles in Dense Sets", "category_id": 6, "category": {"id": 6, "name": "geometry", "display_name": "Geometry", "description": "Euclidean and non-Euclidean geometry, geometric structures.", "slug": "geometry", "order_index": 6, "created_at": "2026-07-31T15:26:25.671Z"}, "status": "open"}
% FIRST_POSTED: 2026-10-01
% Imported from: unsolvedmath_additional_500.tex; UnsolvedMath ID 139
% !TeX program = lualatex
% Compile twice: lualatex 139.tex
% Self-contained: no external bibliography, images, or \input files.
% Numeric section labels and visible numbers are original UnsolvedMath IDs.
\documentclass[11pt,a4paper]{article}
\usepackage[margin=24mm,headheight=14pt]{geometry}
\usepackage{fontspec}
\setmainfont{Latin Modern Roman}
\setsansfont{Latin Modern Sans}
\setmonofont{Latin Modern Mono}
\usepackage{amsmath,amssymb,mathtools}
\usepackage{unicode-math}
\setmathfont{Latin Modern Math}
\usepackage{microtype}
\usepackage{enumitem,xurl,needspace}
\usepackage[dvipsnames]{xcolor}
\usepackage{hyperref}
\hypersetup{unicode=true,colorlinks=true,linkcolor=MidnightBlue,urlcolor=MidnightBlue,
 pdftitle={UnsolvedMath Problem 139},pdfauthor={Source-annotated compilation},bookmarksnumbered=true}
\usepackage{bookmark,tocloft,titlesec,fancyhdr}
\setlength{\cftsecnumwidth}{4.2em}
\cftsetpnumwidth{2.5em}
\cftsetrmarg{4em}
\titleformat{\section}{\Large\bfseries\raggedright}{\thesection}{0.7em}{}
\pagestyle{fancy}\fancyhf{}
\fancyhead[L]{\small\sffamily UnsolvedMath research attempt}
\fancyhead[R]{\small Original catalogue IDs}
\fancyfoot[C]{\thepage}
\setlength{\parindent}{0pt}
\setlength{\parskip}{5pt plus 1pt}
\setlength{\emergencystretch}{3em}
\setcounter{tocdepth}{1}\setcounter{secnumdepth}{1}
\setlist[itemize]{leftmargin=3em,itemsep=4pt,topsep=4pt,parsep=0pt}
\allowdisplaybreaks
\newcommand{\N}{\mathbb{N}}
\newcommand{\Z}{\mathbb{Z}}
\newcommand{\Q}{\mathbb{Q}}
\newcommand{\R}{\mathbb{R}}
\newcommand{\C}{\mathbb{C}}
\newcommand{\F}{\mathbb{F}}
\newcommand{\A}{\mathbb{A}}
\newcommand{\PP}{\mathbb{P}}
\newcommand{\E}{\mathbb{E}}
\newcommand{\eps}{\varepsilon}
\newcommand{\dd}{\,\mathrm d}
\newcommand{\sref}[2]{\hyperlink{source:#1:#2}{[S#2]}}
\newif\ifshowcatalogue
\showcataloguetrue
\newcommand{\cataloguescope}[1]{\ifshowcatalogue
 \paragraph{Statement recorded in the supplied source.}
 #1\fi}
\begin{document}
% UnsolvedMath ID 139; source catalogue code GREEN-051

\setcounter{section}{138}

\section{Large Axis-Parallel Rectangles in Dense Planar Sets}\label{139}

{\small\sffamily UnsolvedMath ID 139 · Geometry}

\subsection{Definitions and mathematical statement}

\paragraph{Shared notation.}Unless a section says otherwise, $\N=\{1,2,\ldots\}$,
$\Z$ is the ring of integers, $\Q,\R,\C$ are the rational, real, and complex
fields, $[N]=\{1,\ldots,\lfloor N\rfloor\}$, and $\log$ is the natural
logarithm. For real functions of a parameter tending to infinity,
$f\ll g$ and $f=O(g)$ mean $|f|\leq Cg$ eventually for some fixed $C$;
$f\gg g$ reverses this comparison; $f=o(g)$ means $f/g\to0$;
$f\sim g$ means $f/g\to1$; and $f\asymp g$ means both order bounds.
A subscript on $O$ or $\ll$ specifies allowed constant dependence.
The expression $n^{o(1)}$ in an upper bound means growth smaller than
$n^\epsilon$ for each fixed $\epsilon>0$. Local definitions govern any
exception, finite-field convention, or use of the same symbol for another
object. An asymptotic claim is not automatically an effective algorithm.



Measure means two-dimensional Lebesgue measure. An axis-parallel rectangle has vertices $(x_1,y_1),(x_1,y_2),(x_2,y_1),(x_2,y_2)$ with $x_1\ne x_2$ and $y_1\ne y_2$; its area is $|x_1-x_2||y_1-y_2|$. Only the four vertices, not the filled rectangle, are required to lie in $A$. The constant $c>0$ must be absolute. \sref{139}{1} \sref{139}{2}

\paragraph{Statement recorded in the supplied source.}
Suppose $A$ is an open subset of $[0, 1]^2$ with measure $\alpha$. Are there four points in $A$ determining an axis-parallel rectangle with area $\geq c\alpha^2$? \sref{139}{1}

\subsection{Short English statement}

Does every open planar set of positive area contain the four corners of an axis-parallel rectangle with area comparable to the square of that density?

\subsection{Research attempt}

\paragraph{Provenance and scope.}
This individual attempt was generated by GPT-6 Astra Ultra on 1 October 2026. The method was to derive explicit partial results and test the first proposed extension adversarially. The original statement and sources below are retained. No claim of mathematical novelty or complete solution is made.

\paragraph{Formal target and context.}
Let $A\subset[0,1]^2$ be open with Lebesgue measure $\alpha>0$. The four vertices, rather than the entire rectangle, must belong to $A$. The target is a rectangle of area at least an absolute constant times $\alpha^2$. The primary problem list supplies this precise density exponent. We attack it by common horizontal fibres, prove an explicit cubic-density lower bound, and construct examples showing why the quadratic exponent is the natural limiting scale.

\paragraph{A common-fibre identity.}
Let $f=1_A$, and for heights $y,z$ put
\[
 R(y,z)=\int_0^1f(x,y)f(x,z)\,dx.
\]
This is the measure of the set of horizontal coordinates at which both rows lie in $A$. Tonelli's theorem and Cauchy--Schwarz give
\[
 \int_0^1\int_0^1 R(y,z)\,dy\,dz
 =\int_0^1\left(\int_0^1 f(x,y)\,dy\right)^2dx\geq\alpha^2.
\]
Merely obtaining a positive value of $R$ does not guarantee a useful area: the rows might be arbitrarily close. We therefore remove a controlled neighbourhood of the diagonal before choosing the pair.

For any $\delta>0$, since $f(x,z)\leq1$, one has
\[
 \iint_{|y-z|<\delta}R(y,z)\,dy\,dz
 \leq\int_x\int_y f(x,y)\left(\int_{|z-y|<\delta}1\,dz\right)dy\,dx
 \leq2\delta\alpha.
\]
Take $\delta=\alpha/4$. The integral over pairs with $|y-z|\geq\delta$ is at least $\alpha^2/2$. As their total parameter area is at most one, there are such pairs with $R(y,z)$ arbitrarily close to or above $\alpha^2/2$. For definiteness choose one with $R(y,z)\geq\alpha^2/4$.

\paragraph{Turning fibre measure into a rectangle.}
A measurable subset $E\subset[0,1]$ of measure $r>0$ has diameter at least $r$: otherwise it would lie in an interval of length less than its measure. Therefore it contains two points separated by more than $r/2$. Apply this to the common fibre of the selected rows. We obtain $x_1,x_2$ with both $(x_i,y)$ and $(x_i,z)$ in $A$, and
\[
 |x_1-x_2|\,|y-z|\geq\frac{\alpha^2}{8}\frac\alpha4
 =\frac{\alpha^3}{32}.
\]
All four points are distinct because both separations are positive. This proves a uniform cubic-density bound. The generous constants avoid any issue about a supremum diameter being unattained. Openness is more than is needed for this argument; measurability and positive measure already suffice.

\paragraph{An explicit quadratic-scale obstruction.}
For an integer $m\geq1$, take the union of the $m$ disjoint open diagonal squares
\[
 A_m=\bigcup_{i=0}^{m-1}(i/m,(i+1)/m)^2.
\]
Its measure is $\alpha_m=1/m$. Suppose all four vertices of an axis-parallel rectangle lie in $A_m$. For any fixed horizontal coordinate $x$, its vertical fibre is exactly the same-index interval containing $x$, or empty if $x$ is a boundary point. Two different index intervals have disjoint fibres. Thus the rectangle's two horizontal coordinates must belong to the same interval, and both vertical coordinates then lie in that interval too. Its width and height are each less than $1/m$, so its area is less than $1/m^2=\alpha_m^2$. Areas approach this value from below by choosing vertices near the square's corners.

Consequently no universal lower bound with a smaller density exponent than two can hold as $\alpha\to0$: a bound $c\alpha^\gamma$ with $\gamma<2$ eventually exceeds $\alpha^2$ on this family. The example does not disprove the desired quadratic bound; it exhibits its sharp scale.

\paragraph{Where the attempted proof loses a factor.}
The common-fibre argument forces vertical separation of order $\alpha$, while it supplies common horizontal measure only of order $\alpha^2$. Their product is cubic. Improving the constant in either estimate cannot change the exponent. To reach the target through this route one needs a joint statement balancing separation against common-fibre diameter, rather than selecting them by two independent lower bounds. The diagonal-square example shows why that balancing is plausible: common fibres are then wider precisely for rows in the same thin block.

\paragraph{Verification and status.}
The diagonal-strip removal uses the actual first moment $\alpha$, not a density-free bound that would lose another power. All integrands are nonnegative, so Tonelli applies without conditional convergence. The example is open in the plane and has exactly the stated area. The checked result is the $\alpha^3/32$ rectangle theorem and the sharpness of exponent two as a possible target. The missing joint geometric inequality remains unproved; the quadratic conclusion is unresolved here.

\subsection{Sources}

{\small

\begin{itemize}

\item[\hypertarget{source:139:1}{S1}] ULAM.AI, \emph{UnsolvedMath}, supplied \texttt{problems.json}, record 139 (GREEN-051), statement and background. The embedded literature-review date is 2026-08-17; that is the archive’s review date, not a fresh certification. Dataset landing page: \url{https://huggingface.co/datasets/ulamai/UnsolvedMath}.

\item[\hypertarget{source:139:2}{S2}] Ben Green, 100 Open Problems (author’s maintained problem list). \url{https://people.maths.ox.ac.uk/greenbj/papers/open-problems.pdf}. The author-hosted document was inspected on 27 September 2026; the archive’s local code is not a problem-number locator in this paper.

\item[\hypertarget{source:139:3}{S3}] GREEN-051 curated problem entry. \url{https://www.unsolvedmath.com/problems/139}. Bibliographic information imported from the supplied record.

\end{itemize}

}

\label{end:139}
\end{document}
